\subsection{Ausgangslage und Problem}\label{sec:ausgangslage_und_problem}

Digitale Plattformen rund ums Kochen gehören heute zum Alltag. Rezept-Communities wie Chefkoch bieten seit Langem Austausch über Gerichte, Bewertungen und Kommentare; daneben hat sich Kochen als eigene Content-Kategorie auf Kurzvideo-Plattformen wie TikTok etabliert.

Parallel dazu gewinnt ein anderes Thema an Dringlichkeit: Lebensmittelverschwendung. Das Nachhaltigkeitsziel 12.3 der Vereinten Nationen fordert, die Lebensmittelverschwendung pro Kopf auf Handels- und Verbraucherebene bis 2030 zu halbieren \parencite[S. 22]{unTransformingOutWorld2030Agenda}. Ein Großteil der Verschwendung wird dabei nicht dem Handel oder der Industrie zugerechnet, sondern der Verbraucherebene und dort vor allem den privaten Haushalten, wo Reste liegen bleiben und schließlich im Müll landen \parencite[S. 105]{vittuariHowReduceConsumer2023}.

Zwischen beiden Entwicklungen zeigt sich eine Lücke. Etablierte Rezept-Communities wie Chefkoch oder die Food-Kategorie auf TikTok organisieren den Austausch über Rezepte, sind aber nicht auf Resteverwertung ausgerichtet: Eine Suche nach einem bestimmten Gericht ist dort leicht möglich, eine Suche nach Rezepten für konkret vorhandene Reste dagegen nicht vorgesehen. Umgekehrt leisten spezialisierte Reste-Matching-Apps wie Restegourmet oder die \ac{BMLEH}-App \enquote{Zu gut für die Tonne!} genau dieses Matching, verfügen aber kaum über Funktionen für Austausch und Community. Hobbyköch:innen, die beides wollen, müssen bislang zwischen mehreren Angeboten wechseln.
